
\section{Experimental Design and Computational Study}
\label{sec:design}

\subsection{Screening and Study Specification}
\label{sec:design:screen}

The specification --- factors, levels, policies, parameters, criteria, primary
decision cost, resolution criterion, descriptor set, selector ladder,
hyperparameters, shift splits and validity rules --- was frozen and committed
before any evaluation run. Everything reported below follows it, with one
recorded exception described in Section~\ref{sec:integrity}.

Before the main campaign, a 16-context screen with ten seeds per cell (640 runs)
evaluated four pre-declared gates on the primary criterion: that at least three
policies be strictly best somewhere with a resolved margin (G1); that at least
25\,\% of contexts have a non-singleton Pareto set (G2); that policy pairs not be
effectively identical on average (G3); and that no single factor explain the
winner map almost entirely (G4). G2, G3 and G4 passed. G1 failed, with two
policies rather than three.

\subsection{Main Factorial Campaign}
\label{sec:design:main}

The main campaign is the $648 \times 4 \times 5$ design: every context under
every policy at every seed, $12{,}960$ runs. It is the basis of every result in
Sections~\ref{sec:validity}--\ref{sec:criteria} and \ref{sec:selection}.

\subsection{Boundary Refinement}
\label{sec:design:boundary}

The main grid samples demand every 600\,veh/h, so it can locate a change in the
preferred policy only to within a 600\,veh/h bracket. Where the resolved winner
changes between adjacent demand levels, a refinement campaign inserts three
intermediate levels at 150\,veh/h spacing inside that bracket, holding all other
factors fixed. Nineteen such transitions were refined, at 1140 runs. A grid
cannot locate a boundary more finely than its own spacing, so the refined result
is reported as a bracket between adjacent sampled levels and never as a
continuous threshold.

\subsection{Sensitivity Experiments}
\label{sec:design:sens}

Two policy parameters had to be chosen and were frozen: the admissibility
tolerance $\epsilon$ of P3 and the risk weight $\lambda$ of P4. A separate
campaign of 1920 runs re-simulates a 96-context subset at
$\epsilon \in \{0.1, 0.3\}$ and $\lambda \in \{0.5, 1.5\}$ to measure how much
the frozen choice matters. The frozen values are not revised in the light of this
experiment; the result is reported as a property of the benchmark
(Section~\ref{sec:robust}).

A second, purely analytical sensitivity varies the preference profile over the
three criteria across three declared normalised weightings, to test whether the
choice of preferred policy depends on the weights.

\subsection{Distribution-Shift Experiments}
\label{sec:design:shift}

Splits O1--O7 are constructed from the main campaign by holding out factor
levels. O8 required a campaign of its own, because no context in the main
factorial contains a disruption anywhere other than on the arterial: 54 contexts
with a north-corridor disruption, 1080 runs.

\paragraph{A planned split that proved infeasible}
The domain-shift split was originally specified as a disruption on the
\emph{south bypass}, which would have been the strongest available test because
the bypass is the reliability-aware policy's preferred corridor. The
corresponding campaign of 2160 runs was generated and executed, and every one of
the 2160 runs failed. The reason is structural rather than computational: the
bypass is a single-lane corridor, so closing a lane on it does not reduce its
capacity, it disconnects it, and the simulator correctly reported that affected
vehicles had no valid route. A lane closure is a capacity reduction only on a
link with more than one lane. The split was replaced by the north-corridor
disruption reported as O8, which is a genuine domain shift --- no training
context contains a disruption on that corridor --- but is a weaker test than the
one originally planned, because the north corridor is not the corridor that any
policy specialises in. The failed campaign is reported in the ledger, the
replacement is not presented as the original design, and the substitution is
counted as a deviation in Section~\ref{sec:integrity}.

\subsection{Campaign Ledger}
\label{sec:design:ledger}

Table~\ref{tab:ledger} reconciles every simulation executed in the study. The
distinction matters because two different totals are defensible and only one of
them supports the results: $12{,}960$ valid runs constitute the main campaign and
every result in Sections~\ref{sec:validity}--\ref{sec:criteria} derives from
them, while $20{,}284$ simulations were attempted across all campaigns, of which
$18{,}124$ succeeded. Neither number is used where the other is meant.

\begin{table}[htbp]
\centering\footnotesize
\caption{Campaign ledger. Every simulation executed in the study, reconciled.
Completion is the fraction of generated vehicles that reached their destination
within the simulated horizon; a run is valid only at completion $\ge 0.98$ with
no teleports.}
\label{tab:ledger}
\setlength{\tabcolsep}{2.6pt}
\begin{tabular}{@{}llrrrrrr@{}}
\toprule
Campaign & Purpose & Runs & Failed & Contexts & Pol. & Seeds & Compl.\\
\midrule
Screen & complementarity screen & 640 & 0 & 16 & 4 & 10 & 1.000\\
Validate & scenario validation & 384 & 0 & 96 & 4 & 1 & 0.973\\
\textbf{Main} & \textbf{main factorial} & \textbf{12{,}960} & \textbf{0} & \textbf{648} & \textbf{4} & \textbf{5} & \textbf{1.000}\\
Boundary & boundary refinement & 1{,}140 & 0 & 19 brack. & 4 & 5 & 1.000\\
Sensitivity & policy parameters & 1{,}920 & 0 & 96 & 2 & 5 & 1.000\\
O8 north & domain shift & 1{,}080 & 0 & 54 & 4 & 5 & 1.000\\
O8 bypass & domain shift, infeasible & 0 & 2{,}160 & --- & 4 & 5 & ---\\
\midrule
\multicolumn{2}{@{}l}{Attempted} & 20{,}284 & & & & &\\
\multicolumn{2}{@{}l}{Succeeded} & 18{,}124 & & & & &\\
\bottomrule
\end{tabular}
\end{table}

\subsection{Experimental Integrity and Deviations}
\label{sec:integrity}

Two deviations from the frozen specification are recorded, and both are reported
here rather than in a supplement because each changes how a reader should weigh
a result.

\paragraph{D-1: the screening gate failed and the campaign was run anyway}
The frozen rule stated that if G1, G3 or G4 failed, the main campaign would not
be run and the study would report a negative result about portfolio design. G1
failed: only two policies, P1 and P3, were strictly best somewhere with a
resolved margin, against a threshold of three. The main campaign was run
regardless. This is a protocol deviation and is not resolved by reinterpreting
the rule.

The reasons were recorded before the campaign was launched and none of them is a
result the campaign might produce. First, the gate's purpose was satisfied: it
exists to prevent a selection study being built on a portfolio with no selection
problem, and a selection problem demonstrably existed, with the preferred policy
reversing between P1 and P3 at resolved margins up to 18.0\,\%. Second, the gate
was written against the primary criterion alone and was blind to the
multi-criteria structure beside it: G2 passed, and all four policies appeared in
Pareto sets, the reliability-aware policy most often of all. Third, the negative
finding was not yet supportable: the screen sampled 16 contexts at extreme factor
levels only and never sampled the demand transition region, so stopping would
have asserted a negative from a small and badly placed fraction of the design.

Four constraints were declared at the same time, before the campaign: that the
study may not claim four-way complementarity on the primary criterion; that it
may not claim P2 or P4 is ever the best journey-time policy unless the full
campaign resolves such a win; that G1's failure must be reported as a finding;
and that a confirmatory gate G1$'$ --- the identical definition evaluated on all
648 contexts --- be declared in advance and its result reported whichever way it
came out. Section~\ref{sec:complementarity} reports that result.

The methodological lesson is recorded as such, because it generalises beyond this
study: a complementarity screen must cover the interior of the factor space and
not only its corners. A resolution-IV fraction at extreme levels is an efficient
design for estimating main effects and a poor one for discovering where a winner
map changes. A gate that is passed only after being overridden was a badly
designed gate.

\paragraph{D-2: the domain-shift split was replaced}
The originally specified O8 was infeasible for the structural reason given in
Section~\ref{sec:design:shift}, and was replaced by a weaker but valid
alternative. The replacement is labelled as such wherever it appears.
